% Generated draft: requires mathematical and source review.
\documentclass[11pt]{article}
\usepackage[margin=25mm]{geometry}
\usepackage{fontspec}
\setmainfont{Latin Modern Roman}
\usepackage{amsmath,amssymb,mathtools,unicode-math}
\setmathfont{Latin Modern Math}
\usepackage{enumitem,needspace,xurl,hyperref}
\setlength{\parindent}{0pt}
\setlength{\parskip}{5pt}
\setlength{\emergencystretch}{3em}
\newcommand{\N}{\mathbb{N}}
\newcommand{\Z}{\mathbb{Z}}
\newcommand{\Q}{\mathbb{Q}}
\newcommand{\R}{\mathbb{R}}
\newcommand{\C}{\mathbb{C}}
\newcommand{\F}{\mathbb{F}}
\newcommand{\A}{\mathbb{A}}
\newcommand{\PP}{\mathbb{P}}
\newcommand{\E}{\mathbb{E}}
\newcommand{\eps}{\varepsilon}
\newcommand{\dd}{\,\mathrm{d}}
\begin{document}
\section*{UnsolvedMath ID: 110}
\textbf{Source problem code:} catalogue entry
\subsection{Definitions and mathematical statement}
The supplied question is an inverse-sieve program: if a small sieve leaves a survivor set of maximal size, characterize its structure. Neither ``small sieve'', ``maximal'', nor the desired notion of structure is defined in the stored statement, and the supplied source assessment explicitly says that a satisfactory precise formulation is unclear.

One candidate model described in that assessment can be defined as follows. For $X\geq2$, choose one excluded class $a_p\in\Z/p\Z$ for each prime $p\leq\sqrt X$, and form
\[
S_X((a_p))=\{n\in\Z:1\leq n\leq\lfloor X\rfloor,
\ n\not\equiv a_p\pmod p\text{ for every prime }p\leq\sqrt X\}.
\]
For this specific model, let
\[
M(X)=\max_{(a_p)_{p\leq\sqrt X}}|S_X((a_p))|.
\]
Exact extremizers satisfy $|S_X|=M(X)$. An asymptotically extremal family, if this relative-size convention is selected, satisfies $|S_X|/M(X)\to1$ as $X\to\infty$. The corresponding inverse problem would classify such choices or survivor sets up to a specified structural equivalence. The source does not select exact versus asymptotic maximality, an error metric, or an equivalence relation, so this model does not by itself yield a uniquely stated classification theorem.

The supplied triage, checked 2026-08-17, reports the Selberg/Rosser--Iwaniec upper-bound-sieve estimate
\[
|S_X((a_p))|\leq(2+o(1))\frac{X}{\log X}
\]
for the one-residue-class-per-prime model. It describes the factor-$2$ extremal phenomenon as connected with a Siegel-zero obstruction and says Green proposes classifying examples arising from that mechanism. No precise Siegel-zero hypotheses, construction, or converse classification are given, so this description is not promoted here into an unconditional assertion that $M(X)\sim2X/\log X$ or that all near-extremizers arise from one mechanism.

The assessment also proposes a different inverse problem at sieve dimension $\kappa=1/2$: it reports that the optimal size is known there and suggests classifying extremizers related to sums of two squares. It does not supply that model's excluded classes, exact dimension definition, optimal-size formula, or equivalence notion. These candidate programs must remain distinct. The source says there is no single settled theorem because its heading intentionally covers a family of inverse problems; the task is first to choose a precise model and then characterize its asymptotically extremal survivors.
\subsection{Short English statement}
Classify the structure of sets that survive a small sieve in maximal or nearly maximal size. The supplied source treats this as a program whose exact formulation is unresolved, with candidate versions involving one excluded residue per small prime, possible Siegel-zero extremizers, or sieve dimension $1/2$.
\subsection{Sources}
\begin{itemize}
\item UnsolvedMath local registry, permanent ID 110, source code catalogue entry, ``Small Sieve Maximal Sets'', Ben Green's 100 Open Problems collection. Dataset attribution: Ulam AI and the UnsolvedMath Contributors; curation and metadata under CC BY 4.0.
\item Ben Green, \emph{100 Open Problems}, Problem 48, January 2026 revision, recorded as checked 2026-08-17, \url{https://people.maths.ox.ac.uk/greenbj/papers/open-problems.pdf}. The supplied triage cites its warning that a good formulation is unclear and its Siegel-zero and sieve-dimension candidate versions.
\item The Selberg/Rosser--Iwaniec bound and the sieve-dimension-$1/2$ progress are supplied through this local assessment; no separate primary theorem citations are given.
\end{itemize}
\end{document}
